% =========================================================
% Existence and Uniqueness Patterns
% =========================================================
\paragraph{Existence and Uniqueness.}

The statement $\exists!\, x\; P(x)$ asserts two things simultaneously:
that some object satisfying $P$ exists, and that it is the \emph{only}
such object. These two obligations are proved separately, in order.

\begin{tcolorbox}[colback=propbox, colframe=propborder, arc=2pt,
  left=6pt, right=6pt, top=4pt, bottom=4pt,
  title={\small\textbf{Template: Existence and Uniqueness ($\exists!\, x\; P(x)$)}},
  fonttitle=\small\bfseries]
\textbf{Step 1 --- Existence.}\\
Define $x := [\text{explicit construction}]$.\\
Verify that $P(x)$ holds. (Check every condition in $P$.)

\medskip
\textbf{Step 2 --- Uniqueness (assume-two-and-compare).}\\
Let $x$ and $y$ be arbitrary objects both satisfying $P$.\\
\emph{Apply the definition of $x$ with $y$ as input; apply the definition
of $y$ with $x$ as input. The two applications collapse.}\\
Conclude $x = y$.

\medskip
\textbf{Step 3 --- Conclude.}\\
By Steps~1 and~2, there exists exactly one object satisfying $P$. \qed
\end{tcolorbox}

\begin{remark}[Existence must be proved before uniqueness]
The uniqueness argument has the form: ``let $x$ and $y$ both satisfy
$P$.'' This assumption is only meaningful if we already know that at
least one object satisfying $P$ exists. Without Step~1, the uniqueness
argument is vacuous: we are asserting something about objects that
might not exist.

This is not pedantry. There are statements where the existence fails
and the uniqueness argument would appear to go through --- but the
whole proof would then be establishing uniqueness for an empty set of
objects, which is trivially true and useless. Existence first is not
just a convention; it is logically necessary.
\end{remark}

\begin{remark}[The key move in assume-two-and-compare]
The assume-two-and-compare argument has a characteristic structure that
appears across mathematics. The critical move is this: if $x$ and $y$
both satisfy the defining property $P$, apply $x$'s property with $y$
as the argument, and simultaneously apply $y$'s property with $x$ as
the argument. The two applications collapse to $x = y$.

Here is the pattern in the group identity example: if $e$ satisfies
the identity axiom, then $e \cdot e' = e'$ (using $e$ as identity with
argument $e'$). If $e'$ satisfies the identity axiom, then $e \cdot e' = e$
(using $e'$ as identity with argument $e$). Combining: $e' = e$.

The same move works for: uniqueness of inverses, uniqueness of the
zero vector, uniqueness of limits, uniqueness of the supremum.
In every case, the two assumed objects are ``played against each other'':
each one's defining property is applied with the other as input.
\end{remark}

\begin{example}[Uniqueness of the group identity]
\textit{Proof.}
Let $G$ be a group. Suppose $e$ and $e'$ both satisfy the identity axiom:
\[
e \cdot g = g \cdot e = g \quad \text{for all } g \in G,
\]
\[
e' \cdot g = g \cdot e' = g \quad \text{for all } g \in G.
\]

Apply $e'$'s identity property with $g = e$:
\[
e' \cdot e = e.
\]
Apply $e$'s identity property with $g = e'$:
\[
e \cdot e' = e'.
\]
But $e' \cdot e = e \cdot e'$ (both equal the product $e \cdot e'$ in $G$
by commutativity --- actually, both equal $e \cdot e'$ directly).
Actually, more directly: $e = e' \cdot e = e'$ where the first step
uses $e'$ as identity and the second uses $e$ as identity.
Hence $e = e'$. \qed

\medskip\noindent
The single line is: $e = e \cdot e' = e'$.
First equality: $e$ acts as identity on $e'$.
Second equality: $e'$ acts as identity on $e$.
Both identities applied to the same product; both sides collapse to
the same expression; $e = e'$.
\end{example}

\begin{remark}[Why proving uniqueness creates a proof tool]
Every uniqueness theorem, once proved, creates a new application of
the satisfy-and-cite tactic. This is the connection that most students
miss.

Proving uniqueness says: there is exactly one object satisfying $P$.
This means: if you can ever show that some element $a$ satisfies $P$,
you automatically get $a = b$, where $b$ is the canonical name for the
unique object satisfying $P$.

The group inverse provides the clearest example. Once you have proved
that inverses are unique, the pattern for proving $(ab)^{-1} = b^{-1}a^{-1}$
is immediate: show $b^{-1}a^{-1}$ satisfies the defining property of
the inverse of $ab$, then cite uniqueness. You never need to set up
another assume-two-and-compare argument. The uniqueness theorem absorbs
all future comparisons.

This is why uniqueness results are not just facts --- they are
\emph{tools}. Prove them early; use them everywhere afterward.
\end{remark}
